%%%%%%%%%%%%%%%%%%%%%%%%%%%%%%%%%%%%%%%%%%%%%%%%%
% EXERCICE
% PGCD et PPCM
\begin{exercise}[points=2]
    Elisa et Angelina font du sport dans un champ. Elisa marche et fait le tour du champ en \SI{12}{min}. Angelina, elle, s’est mise à la course et fait le tour du champ en \SI{5}{min}. Si elles démarrent en même temps d’un point donné à 13 h 45, à quelle heure vont-elles repasser à ce même endroit en même temps ?\\
    Laisse ta démarche visible.
\end{exercise}
% SOLUTION
\begin{solution}

\end{solution}

\vspace{7cm}
%%%%%%%%%%%%%%%%%%%%%%%%%%%%%%%%%%%%%%%%%%%%%%%%%
%%%%%%%%%%%%%%%%%%%%%%%%%%%%%%%%%%%%%%%%%%%%%%%%%
% EXERCICE
% Division euclidienne
% \begin{exercise}
%     Magali possède 93 figurines de princesses. Elle souhaite les ranger dans des boîtes qui peuvent en contenir 8. Combien de boîtes devront être utilisées pour ranger toutes les princesses de Magali ?
% \end{exercise}
% % SOLUTION
% \begin{solution}

% \end{solution}
%%%%%%%%%%%%%%%%%%%%%%%%%%%%%%%%%%%%%%%%%%%%%%%%%
%%%%%%%%%%%%%%%%%%%%%%%%%%%%%%%%%%%%%%%%%%%%%%%%%

\newpage
\begin{exercise}[points=3]
    Martin possède un coffre de rangement ayant la forme d’un parallélépipède rectangle dont les dimensions intérieures sont \SI{32}{cm}, \SI{40}{cm} et \SI{56}{cm}. \\
Il veut y ranger des cubes aussi grands que possible dont l’arête est mesurée par un nombre entier de centimètres.
\begin{enumerate}
    \item Détermine la mesure de l'arête du plus grand cube possible.
    \item Combien de cubes pourra-t-il mettre dans son coffre de rangement?
\end{enumerate}
Laisse tes étapes visibles.
\end{exercise}
%SOLUTION
\begin{solution}
    
\end{solution}

\vspace{7cm}
%%%%%%%%%%%%%%%%%%%%%%%%%%%%%%%%%%%%%%%%%%%%%%%%%
%%%%%%%%%%%%%%%%%%%%%%%%%%%%%%%%%%%%%%%%%%%%%%%%%
\begin{exercise}[points=3]
    Valérie a reçu des résultats à un test d'EDM. Elle a obtenu la note de $50/70$. Caleb quant à lui a obtenu la note de $60/80$. Qui a obtenu le meilleur résultat ?\\
    Laisse ta démarche visible.
\end{exercise}
% SOLUTION
\begin{solution}

\end{solution}

%%%%%%%%%%%%%%%%%%%%%%%%%%%%%%%%%%%%%%%%%%%%%%%%%
%%%%%%%%%%%%%%%%%%%%%%%%%%%%%%%%%%%%%%%%%%%%%%%%%

\vspace{4cm}
%%%%%%%%%%%%%%%%%%%%%%%%%%%%%%%%%%%%%%%%%%%%%%%%%
%%%%%%%%%%%%%%%%%%%%%%%%%%%%%%%%%%%%%%%%%%%%%%%%%
\begin{exercise}[points=2]
    Un c\oe ur humain bat en moyenne $ 1,15 \cdot 10^5$ fois par jour. Combien de fois un cœur humain bat-il sur une année? \\
    Laisse ta démarche visible et écris ton résultat en notation scientifique.
\end{exercise}

\newpage
% Notation scientifique
\begin{exercise}[points=3]
    À l'aide du tableau ci-dessous, réponds aux questions.
    \begin{center}
        \begin{tblr}{hlines, vlines,row{1}={m,font=\bfseries}, columns={c, 4cm}, rowsep=1ex}
            Planète & Distance par rapport au soleil (km) & Notation scientifique\\
            %Mercure & \num{58000000}       \\
            Vénus   & \num{108 190 000} &   \\
            Terre   & \num{149 569 000} &   \\
            Mars    & \num{227 940 000} &   \\
            Jupiter & \num{778 000 000} &   \\
            %Saturne & \num{1 427 000 000}  \\
            %Uranus  & \num{2 871 000 000}  \\
            %Neptune & \num{4 497 000 000}
        \end{tblr}
    \end{center}
    \begin{enumerate}[itemsep=3em]
        \item Convertis ces distances en notation scientifique.
        \item Laquelle de ces planètes se trouve à environ $2\cdot 10^8$ km du soleil ?
        \item Dans le système solaire, on peut retrouver des milliers d'astéroïdes. Ceux-ci sont beaucoup plus petits que les planètes. Il y a une immense ceinture d'astéroïdes à environ $6\cdot 10^8$ km du soleil. Entre quelles planètes se trouvent cette ceinture ? 
    \end{enumerate}
\end{exercise}
% SOLUTION
\begin{solution}

\end{solution}
%%%%%%%%%%%%%%%%%%%%%%%%%%%%%%%%%%%%%%%%%%%%%%%%%
%%%%%%%%%%%%%%%%%%%%%%%%%%%%%%%%%%%%%%%%%%%%%%%%%
<<<<<<< Updated upstream
\vspace{2cm}


=======
%calcul un peu vache donc on pourrait le mettre sur peu pour que ca n'influence pas trop. 
%Ce n'est qu'une proposition, 1pt pour multiplier H par 2, 1pt pour diviser O2 par deux et 1pt pour l'addition.
%On peut rajouter des infos pour faciliter ou laisser tomber si trop compliqué :)
\newpage
\begin{exercise}[points=3]
Sachant que $6,023.10^{23}$ molécules d'hydrogene (H) pèsents 1 gramme
et que $6,023.10^{23}$ molécules de dioxygene ($O_2$) pèsent 32 grammes, 
combien de molécules y aurait-t-il dans 18 grammes d'$h_2O$ ?
    
\end{exercise}
\vspace*{4cm}
%%%%%%%%%%%%%%%%%%%%%%%%%%%%%%%%%%%%%%%%%%%%%%%%
% 1 pt par réponse
\begin{exercise}[points=3]
    On souhaite paver une terrasse de 4,20m sur 3,60m avec des dalles carrées.\\
    a) Détermine la longueur du coté des plus grandes dalles qui peuvent être utilisées si on ne veut pas faire de découpes.\\
    b) Détermine le nombre de dalles qui sera nécessaire pour réaliser ce pavement.\\
    c) Le propriétaire trouve ces dalles trop grandes et en choisit d'autres dont le côté est réduit de moitié.
    Détermine le nombre de dalles qu'il devra commander.
\end{exercise}
>>>>>>> Stashed changes
